%%%------------------------------------
%%%     DIPLOMARBEIT: Hauptdatei
%%%     
%%%     Stefan Grünwald (2000 - 2007)
%%%     Michael Fladischer (2010)
%%%
%%%     Das Dokument lässt sich am besten mit "latexmk" bauen.
%%%------------------------------------

% Dokumentenklasse%{{{
\documentclass[
    %DIV12,                                 % Satzspiegelparameter 12
    %BCOR8mm,                               % Bindekorrektur 8mm
    %headinclude,                           % Berücksichtigung der Kopfzeile
                                            % bei der Ermittlung des Satzspiegels
    12pt,                                   % Schriftgrösse 12pt
    a4paper,                                % DIN A4
    oneside,                                % einseitig
    captions=oneline,                       % 
    ngerman,
    draft=true,                             % Muss entfernt werden, bevor Arbeit abgegeben wird!
    bibliography=totoc,                     % Literaturverzeichniss automatisch im TOC einbinden
    listof=totoc,                           % Listen automatisch in TOC einbinden
    index=totoc]                            % Index automatisch in TOC einbinden
    {scrbook}                               % KOMA-Script Klasse 'scrbook' statt 'book'
%}}}

% Generische Packages%{{{
\usepackage[ngerman]{babel}
\usepackage[ngerman]{translator}
\usepackage[utf8]{inputenc}                 % Laden der UTF8-Eingabekodierung
                                            % ACHTUNG: Dateien müssen UTF8 codiert sein!
\usepackage[T1]{fontenc}                    % EC codierte Fonts
\usepackage{ae}                             % PDF encoding Info
\usepackage{palatino}                       % Brotschrift: Palationo
\usepackage{helvet}                         % Serifenlose Schrift: Helvetica
\usepackage[left]{eurosym}                  % EURO links
\usepackage{paralist}                       % Beeinflussung von Listen
\usepackage{tabularx}                       % Austreiben auf eine bestimmte Breite
\usepackage{longtable}                      % Mehrseitige Tabellen
\usepackage{ltxtable}                       % 'longtable' und 'tabularx'
\usepackage{array}                          % Zeilenumbruch in 'p'-Spalten
\usepackage{booktabs}                       % \toprule, \midrule und \bottomrule
\usepackage{dcolumn}                        %
\usepackage{blindtext}                      % Erzeugt Blindtext welcher beim Prüfen des Layouts behilflich sein kann
\usepackage{siunitx}
\usepackage{color}
\usepackage{xcolor}
% http://www.ctan.org/tex-archive/macros/latex/contrib/wasysym/
\usepackage{wasysym}
% http://www.ctan.org/tex-archive/macros/latex/contrib/url/
\usepackage[hyphens]{url}
%}}}

% Formatierung für Listings%{{{
\usepackage{listings}
\lstloadlanguages{HTML,XML}
\lstset{float,
        numbers=none, 
        keywordstyle=\bfseries\emph, 
        identifierstyle=\ttfamily, 
        emphstyle=\color{red}\bfseries\emph,
        columns=fixed,
        prebreak=\raisebox{0ex}[0ex][0ex]{\ensuremath{\hookleftarrow}},
        breaklines=true,
        firstnumber=auto,
        numbers=left, 
        numberstyle=\tiny, 
        stepnumber=2, 
        numbersep=5pt,
        xleftmargin=12pt,
        escapeinside={@}{@},
        basicstyle=\ttfamily}%}}}

% Lokale Einstellungen%{{{
% vim: set tw=160:

\def\campusType{DIPLOM-ARBEIT}
\def\campusTitle{To be defined}
\def\campusAuthor{Michael Fladischer}
\def\campusPersKnz{0810319004}
\def\campusCity{Graz}
\def\campusDate{\today}
\def\campusSubject{}
\def\campusKeywords{PGP,x509,Firefox,XML-DSig}

%}}}

% Weiche für PDF/DVI%{{{
\usepackage{ifpdf}
% für Encapsulated-Postscript und PDFTeX Grafiken
\ifpdf
    \usepackage[pdftex]{graphicx}
    \graphicspath{{images/pdf/}}
    \pdfcompresslevel=9
\else
    \usepackage[dvips]{graphicx}
    \graphicspath{{images/dvi/}}
\fi%}}}

% TIKZ - Coole Grafiken direkt in TeX zeichnen :-)%{{{
\usepackage{tikz}
\usetikzlibrary{mindmap,trees,arrows,shapes,positioning,shadows,calc}
%}}}

% Überschriften und Zeilenabstand%{{{
\usepackage{layout}

% Überschriften formatieren
\usepackage[font=small,format=hang]{caption}

% Überschriftenformat für Listings
\DeclareCaptionFont{white}{\color{white}}
\DeclareCaptionFormat{listing}{\colorbox{gray}{\parbox{\textwidth}{#1#2#3}} }
\captionsetup[lstlisting]{format=listing,labelfont=white,textfont=white}

\renewcommand{\floatpagefraction}{.8}       % ,float page`, wenn diese mindestens
                                            % 80% der Seite einnehmen
                                            % default: .5
\renewcommand{\textfraction}{.1}            % neben den Abbildungen, die ,h`ere, ,t`op und
                                            % ,b`ottom plaziert werden, noch min. 10% Text
                                            % default: .2
\renewcommand{\topfraction}{.9}             % max. Anteil der Seite, die f\"{u}r Abbildungen,
                                            % die ,t`op plaziert werden, verwendet wird
                                            % default: .7
\renewcommand{\bottomfraction}{.6}          % max. Anteil der Seite, die f\"{u}r Abbildungen,
                                            % die ,b`ottom plaziert werden, verwendet wird
                                            % default: .3
%}}}

% Silbentrennung%{{{
% vim: set tw=160:

\hyphenation{Ge-sch\"{a}fts-pro-zes-sen
             Ge-sch\"{a}fts-pro-zess-mo-dell
}


%}}}

% Index-Erstellung%{{{
\usepackage{makeidx}
\makeindex
%}}}

% Smarte Referenzen%{{{
\usepackage[german]{fancyref}
% Listings (lst:my-unigue-label)
\newcommand*{\fancyreflstlabelprefix}{lst}
% Zeilen in Listings (lin:my-unique-label)
\newcommand*{\fancyreflinlabelprefix}{lin}
\fancyrefaddcaptions{german}{
    \newcommand*{\Freflstname}{\lstlistingname}
    \newcommand*{\freflstname}{\lstlistingname}
    \newcommand*{\Freflinname}{Zeile}
    \newcommand*{\freflinname}{Zeile}
}
\frefformat{vario}{\fancyreflstlabelprefix}{\Freflstname\fancyrefdefaultspacing#1#3}
\frefformat{vario}{\fancyreflinlabelprefix}{\Freflinname\fancyrefdefaultspacing#1#3}
%}}}

% PDF-Einstellungen%{{{
% http://www.ctan.org/tex-archive/macros/latex/contrib/hyperref/
\usepackage[a4paper,                        % DIN A4
            plainpages=false,
            pdfpagelabels,
            pdftitle={\campusTitle},
            pdfauthor={\campusAuthor},
            pdfsubject={\campusType, \campusDate, \campusSubject},
            pdfkeywords={\campusKeywords},
            bookmarks=true,                 % Lesezeichen im 'pdf' analog zum Inhaltsverzeichnis,
            bookmarksopen=true,             % Bookmarks anzeigen...
            bookmarksnumbered=true,         % ...und nummerieren
            hyperindex=true,                % Indexeinträge als Hyperlinks
            colorlinks,                     % Farbige Links
            linkcolor=black,
            urlcolor=black,
            citecolor=black,
            pdfborder=0,
            unicode=true,
            breaklinks
]{hyperref}

% Allow links to have line-breaks even when processed by dvips/ps2pdf.
% http://www.ctan.org/tex-archive/macros/latex/contrib/breakurl/
\usepackage{breakurl}
%}}}

% Glossaries%{{{
\usepackage[toc,acronym,style=long]{glossaries} 
\loadglsentries{misc/glossary}
\makeglossaries
%}}}

% Benutzerdefinierte Befehle%{{{
\newcommand{\xmlelem}[1]{\index{#1}\texttt{<#1/>}}
\newcommand{\xmlattr}[2]{\index{#1!#2}\textit{#2}}
\newcommand{\yes}{$\circledast$}
\newcommand{\no}{$\circleddash$}
\newcommand{\optional}{$\circledcirc$}
%\newcommand{\yes}{\rlap{\Cross}\Square}
%\newcommand{\no}{\Square}
%\newcommand{\optional}{\rlap{\FilledSmallCircle}\Square}
%\newcommand{\yes}{\rlap{\FilledSmallCircle}\Circle}
%\newcommand{\no}{\Circle}
%\newcommand{\optional}{\rlap{\SmallCircle}\Circle}
%\newcommand{\todo}{\marginpar{\textbf{TODO}}}
\renewcommand*{\glspostdescription}{}

%}}}

% Layout-Einstellungen%{{{

% Numerierungstiefe der Kapitel
\setcounter{secnumdepth}{3}%

% Numerierungstiefe des Inhaltsverzeichnisses
\setcounter{tocdepth}{3}%

% Abstand vor und nach Gleitobjekten im Text [h]
%\setlength{\intextsep}{12pt plus 2pt minus 2pt}    % default: {12pt plus 2pt minus 2pt}

% Abstand zwischen Gleitobjekten oben oder unten und Text [tb]
\setlength{\textfloatsep}{28pt plus 2pt minus 4pt}  % default: {20pt plus 2pt minus 4pt}

%\setlength{\abovecaptionskip}{0pt}

%%%---- PARAGRAPH ---------------------
\makeatletter
\renewcommand{\paragraph}{%
  \@startsection {paragraph}{4}{\z@}%
  {-3.25ex plus-1ex minus-.2ex}{1sp}%
  {\raggedsection \normalfont \sectfont \nobreak \size@paragraph \nobreak } % was: \reset@font \normalsize\bf
} \makeatother

%%%---- SUBPARAGRAPH ------------------
\makeatletter
\renewcommand{\subparagraph}{%
  \@startsection {subparagraph}{5}{\parindent}%
  {3.25ex \@plus 1ex \@minus .2ex}{-1em}%
  {\raggedsection \normalfont \sectfont \nobreak \normalsize\bf \nobreak } % was: \size@subparagraph
} \makeatother

\clubpenalty = 10000
\widowpenalty = 10000
% \displaywidowpenalty = 10000

%\hyphenpenalty=9999                     % verhindert häufige Silbentrennung
%\lefthyphenmin=3                        % bei Silbentrennung bleiben min. 3 Zeichen
                                         % in der oberen Zeile stehen
%\finalhyphendemerits=5000%
%\doublehyphendemerits=10000%
%\sloppy                                 % sollte unbedingt vermieden werden ;-)

%%%------------------------------------
%%%---- SEITENDEKLARATION -------------
%%%------------------------------------
\setlength{\textheight}{660pt}%
\setlength{\voffset}{40pt}%
\setlength{\hoffset}{20pt}%
%\setlength{\footskip}{20pt}%

%%%------------------------------------
%%%---- KOPF-, FUSSZEILE --------------
%%%------------------------------------
\usepackage{fancyhdr}                   % Laden des Pakets zum Layout
%\pagestyle{fancy}                       % der Kopf- und Fusszeilen
\renewcommand{\chaptermark}[1]{\markboth{\thechapter.\ #1}{}}
%\renewcommand{\sectionmark}[1]{\markright{#1}{}}
%\renewcommand{\headrulewidth}{0.4pt}    % Linie Kopfzeile
%\renewcommand{\footrulewidth}{0.4pt}   % Linie Fusszeile
%\lhead{\sffamily\small\leftmark}        % LeftHeader --> Kapitelname
%\chead{}                                % CenteredHeader
%\rhead{\sffamily\small\thepage}         % RightHeader --> Seitennummer
%\lfoot{\sffamily\small}                 % LeftFooter --> Name
%\cfoot{}                                % CenteredFooter
%\rfoot{\sffamily\small}                 % RightFooter --> Titel
%%%------------------------------------
%}}}

% Haupt-Dokument%{{{

\begin{document}

%%%------------------------------------
%%%---- VORSPANN ----------------------
%%%------------------------------------
\frontmatter                                % Titel, Vorwort, Inhaltsverzeichnis

%%%---- TITELSEITE --------------------
\begin{titlepage}
\begin{sf}
\begin{center}
\huge \campusType \\

\vspace{2cm}%

\campusTitle \\
\vspace{3cm}%

\begin{normalsize}

ausgeführt am

\begin{figure} [h!]
        \centerline{\includegraphics[width=0.3\textwidth]{campus02}}
\end{figure}

%\vspace{4.5cm}%

Studiengang\\ 
Informationstechnologien und IT-Marketing\\

\vspace{3.5cm}%

%\begin{tabularx}{\linewidth}{ll}
Von: \campusAuthor \\
Pers. Kennz. \campusPersKnz
%\end{tabularx}

\vspace{3.5cm}%
\end{normalsize}
\end{center}

\begin{tabularx}{0.9\linewidth}{Xc}
\campusCity, am \campusDate & ............................. \\
         & \campusAuthor\\
\end{tabularx}

\end{sf}
\end{titlepage}

\newpage

% Seitennummerierung auf Römisch setzen und bei Seite 1 beginnen
\pagenumbering{roman}
\setcounter{page}{1}

%%%---- EO TITELSEITE -----------------

\setlength{\parindent}{0em}

%%%---- EHRENWORT ---------------------
%%%%%%%%%%%%%%%%%%%%%%%%%%%%%%%%%%%%%%%
%%%     EHRENW\"{O}RTLICHE ERKL\"{A}RUNG
%%%     2007/10/17
%%%%%%%%%%%%%%%%%%%%%%%%%%%%%%%%%%%%%%%

\chapter*{Ehrenw\"{o}rtliche Erkl\"{a}rung}

Ich erkl\"{a}re ehrenw\"{o}rtlich, dass ich die vorliegende
Arbeit selbstst\"{a}ndig und ohne fremde Hilfe verfasst, andere als
die angegebenen Quellen nicht ben\"{u}tzt und die benutzten Quellen
w\"{o}rtlich zitiert sowie inhaltlich entnommene Stellen als solche
kenntlich gemacht habe. \\

\vspace{1 cm} \hfill
\begin{tabular*}{4.4 cm}{c}
........................................ \\
Unterschrift \\
\end{tabular*}

%%%---- DANK --------------------------
%%%%%%%%%%%%%%%%%%%%%%%%%%%%%%%%%%%%%%%
%%%     DANKSAGUNG
%%%     2007/10/17
%%%%%%%%%%%%%%%%%%%%%%%%%%%%%%%%%%%%%%%

\chapter*{Danksagung}

.....

\vspace{1.5cm}

\begin{center}
\raggedleft{\campusAuthor}
\end{center}

\vspace{1.5cm}

\campusCity, am \campusDate


%%%---- KURZFASSUNG -------------------
% vim: set tw=160:

% 2010 Peter A. Menzel

\noindent
\begin{minipage}{\textwidth}
%%% minipage um Kapitel Kurzfassung UND Kapitel Abstract
%%% auf eine Seite zu binden
\chapter*{Kurzfassung}
Die zunehmende Verbreitung von Websites, auf denen es möglich ist, auch als Benutzer Inhalte zu veröffentlichen, macht es immer wichtiger, diese Inhalte auch
eindeutig, über die Grenzen einer einzelnen \glstext*{website} hinweg, einem Autor eindeutig zuordnen zu können. Um diese Zuordnung zuverlässig herstellen zu
können, bieten sich die drei kryptographischen Eigenschaften Integrität, Authentizität und Verbindlichkeit in der Form von digitalen Signaturen an. Dazu
betrachtet diese Arbeit Anforderungen, die an digitale Signaturen in einer \glstext*{website} gestellt werden. Es werden Ansätze zur Verbindung von
\gls{website} und digitaler Signatur erläutert und die Verfahren \gls{openpgp} und \gls{x509} auf ihre Tauglichkeit für diesen Einsatzzweck gegenübergestellt.
Als Resultat werden aus diesen Informationen mehrere Anforderungen für ein, noch zu entwickelndes, \glstext{firefox-addon} für \glstext*{xml-dsig} abgeleitet.

\chapter*{Abstract}
With the proliferation of websites that offer their users to publish their own content, it is becoming more and more important, that content can be assigned an
author even across the boundaries of a single website. To establish such a reliable mapping it is crucial to be able to apply the three cryptographic
properties of integrity, authenticity and non-repudiability in the form of digital signatures to any content of a website. This work researches requirements
imposed by digital signatures in a website. It identifies approaches to embedding signatures within a website. For this purpose two digital signature systems,
\glstext*{openpgp} and \glstext*{x509}, are evaluated on their suitability in a website scenario. As a result, this information is used to derive requirements
for a, yet to be developed, Firefox add-on that incorporates \glstext*{xml-dsig} to validate digital signatures in \glsplural*{website}.

\end{minipage}


% Gleichheitsgrundsatz
% vim: set tw=160:

\chapter*{Gleicheitsgrundsatz}
Aus Gründen der Lesbarkeit wurde in dieser Arbeit darauf verzichtet, geschlechtsspezifische Formulierungen zu verwenden. Jedoch möchte der Author ausdrücklich
festhalten, dass die bei Personen verwendeten femininen oder maskulinen Formen für beide Geschlechter zu verstehen sind.


%%%---- INHALTSVERZEICHNIS ------------
\tableofcontents

%%%------------------------------------
%%%---- HAUPTTEIL ---------------------
%%%------------------------------------
\setlength{\parskip}{3ex plus0.5ex minus0.5ex}

%%%---- SEITENLAYOUT ------------------
%\newpage \layout

\mainmatter
% Seitennummerierung auf Arabisch setzen und bei Seite 1 beginnen
\pagenumbering{arabic}
\setcounter{page}{1}

% Laden des eigentlichen Inhalts
% vim: set tw=160:

\chapter{Einleitung}
\index{Einleitung}%
\label{chap:Einleitung}

\section{Ausgangssituation}
\index{Ausgangssituation}%
\label{sec:Einleitung:ausgangssituation}
Kryptographische Merkmale wie Vertraulichkeit, Authentizität und Integrität von digitalen Inhalten gewinnen in einem wachsenden globalen Netzwerk immer
stärker an Bedeutung. So lag die Anzahl von \gls{ssl}-fähigen Webservern bei \num{10.8e6} und den darauf eingesetzten gültigen Zertifikaten im Juli 2010 bei
\num{4.3e6} \cite{ssliverse:eff}.

\gls{ssl}/\gls{tls} dient jedoch nur der Vertraulichkeit sowie der Integrität der Gesamtheit aller übertragenen Daten und der Authentizität zwischen den Peers einer unidirektionalen
Kommunikation, wie sie zum Beispiel beim \gls{https} zur Anwendung kommt \cite{kits}. Aktive (umleiten der Verbindung) und passive Angriffe (Lesen der übertragenen Daten)
lassen sich dadurch verhindern. Eine Prüfung, ob die der gesamte übertragen Inhalt, oder Teile dessen, authentischen Ursprungs sind, ist damit jedoch nicht
möglich.

In anderen Bereichen der Kommunikation wie E-Mail und \gls{im} gehören Verfahren zum Nachweis der Identität des Absenders und der Integrität der Nachricht
bereits zum Leistungsumfang vieler Anwendungsprogramme \cite{2719799020071101}.\\
Die Verfahren, welche dabei die größte Verbreitung finden, sind \gls{pgp} und \gls{smime}. Sie bieten die Möglichkeit, digitale Signaturen auf Basis von asymmetrischen
Verschlüsselungsverfahren zu erstellen \cite{kits}. Dabei wird vom Verfasser einer Nachricht diese mit seiner digitalen Signatur versehen. Der Empfänger ist damit
in der Lage, die in der Nachricht signierten Teile zu verifizieren.

Mit dem Aufkommen des Web 2.0 entstand auf vielen \glspl{website} für die Besucher die Möglichkeit, eigene Inhalte einzubinden oder bestehende Inhalte weiter zu
verwenden. Der Betreiber einer \gls{website} wurde dadurch vom alleinigen Bereitsteller zu einem Verwalter von Informationen. Für Besucher solcher Seiten mit
gemischten redaktionellen und benutzergenerierten Inhalten existiert noch keine Möglichkeit, diese Trennung auch kryptografisch sicher nachvollziehen zu können.
Das Vertrauen eines Besuchers in die Identität des Urhebers und die Integrität der Inhalte kann deshalb nur über die optische Darstellung gebildet werden. Eine
optische Trennung von Inhalten bietet jedoch noch keine gesicherte Information über Authentizität und Integrität. Hier bieten sich digitale Signaturen  an, um
genau dieses Bedürfnis nach kryptografisch sicherer Auszeichnung von Inhalten einer \gls{website} erfüllen.

\section{Aufgabenstellung}
\index{Aufgabenstellung}%
\label{sec:Einleitung:aufgabenstellung}
Die primäre Aufgabenstellung dieser Arbeit besteht darin, zu bestimmen, welche Anforderungen an eine Integration von digitalen Signaturen in die
\gls{xhtml}-Dokumente einer \gls{website} gestellt werden. Folgende Aspekte sollen deshalb in dieser Arbeit erläutert werden:
\begin{itemize}
    \item Wie können digitale Signaturen in \gls{xhtml} integriert werden, ohne dass der entsprechende W3C-Standard \cite{xhtml:w3c} erweitert oder angepasst
    werden muss?
    \item Welche Transformationen von \gls{xml-dsig} eignen sich zur wahlfreien Selektion von Bereichen innerhalb eines \gls{xhtml}-Dokuments?
    \item Wie können die öffentlichen Schlüssel zu den Signaturen verteilt werden? Welche der bestehenden Technologien eignen sich dafür am besten?
\end{itemize}

\section{Zielsetzung der Arbeit}
\index{Ziele}
\label{sec:Einleitung:ziele}
Diese Arbeit wird die Eignung der vorhandenen Signaturverfahren für \gls{xml}, sowie die damit einhergehenden Möglichkeiten zur Schlüsselverteilung
behandeln. Um den praktischen Nutzen eines solchen Verfahrens zu demonstrieren soll auch aufgezeigt werden, welche Arten von Angriffen auf Inhalte durch
digitale Signaturen in \gls{xhtml} verhindert werden können.

Nicht behandelt wird die Erstellung eines neuen Signaturverfahrens, eines neuen Signaturmodells für \gls{xml} oder die Erweiterung des bestehenden
\gls{xhtml}-Standards. Auch ist eine Implementierung eines Browser-Plugins nicht im Rahmen dieser Arbeit vorgesehen. Die \gls{html} wird nicht behandelt, da es
sich dabei nicht um ein mit \gls{xml} kompatibles Format handelt.

\section{Vorgehen und Methodik}
\index{Methodik}
\label{sec:Einleitung:methodik}


\section{Aufbau der Arbeit}
\index{Aufbau der Arbeit}
\label{sec:Einleitung:aufbau}
Eine kurze Einführung in \gls{xml-dsig} und verwandte Verfahren zur Signierung von Inhalten, danach Überleitung zur konkreten Arbeitsweise von \gls{xml-dsig}.
Vorstellung der Möglichkeiten zur Verbindung zwischen \gls{xml-dsig} und \gls{xhtml}. Evaluierung und Auswahl des am besten geeigneten Mechanismus, bewertet an
Praktikabilität und Sicherheit. Überblick über die von \gls{xml-dsig} unterstützten Signaturverfahren. Evaluierung der Verfahren und Auswahl eines Verfahrens
anhand der Eignung für den Einsatz in einem \gls{webbrowser}.


% vim: set tw=160:

\chapter{Grundlagen und Definitionen}
\label{chap:GrundlagenDefinitionen}

Zum Verständnis der folgenden Kapitel wird ein grundlegendes Verständnis von \gls{xml}, \gls{xhtml} und asymmetrischer Kryptographie vorausgesetzt.
Themenbereiche wie digitale Signaturen und für diese Arbeit relevanten Techniken, die vor allem in Verbindung mit \glspl{website} vorkommen, werden hier kurz
dargestellt. Dazu gehören vor allem \gls{openpgp}, \gls{x509} und \gls{xml-dsig}.

\section{\glsentrytext{website}}
\index{Website}
\label{sec:GrundlagenDefinitionen:Website}
Eine \gls{website} repräsentiert eine Menge von Dateien, die unter dem gemeinsamen Namensraum einer, oder auch mehrerer Domains mit dem selben Bezug, abrufbar
sind. Zu diesen Dateien gehören vorrangig \gls{html}- und \gls{xhtml}-Dokumente, \gls{css} und \gls{javascript}, sowie auch komplementäre Bild-, Ton- und
Videodateien.

\section{\glsentrytext{webseite}}
\index{Webseite}
\label{sec:GrundlagenDefinitionen:Webseiten}
Der Begriff einer \gls{webseite} bezeichnet ein einzelnes \gls{html}- oder \gls{xhtml}-Dokument inklusive aller von ihm zur Darstellung referenzierten
Inhalte wie \gls{css}-Dateien und \gls{javascript}, sowie auch komplementäre Bild-, Ton- und Videodateien. Inhalte, auf die mittels \gls{html}-Links, wie sie
von \xmlelem{a}-Elementen erzeugt werden, verwiesen wird, gehören nicht zum Umfang einer Webseite, da sie nicht direkt in die Darstellung in einem
\gls{webbrowser} einfließen.

\section{Bereiche in \glsentrytext{xhtml}-Dokumenten}
\label{sec:GrundlagenDefinitionen:BereicheXHTML}
Der Dokumententyp \gls{xhtml} ist eine Anwendungsform von \gls{xml} auf das \gls{sgml}-Format \gls{html}\cite{xhtml:w3c} und kann dadurch in seiner Form als
\gls{dom}-Baum betrachtet werden \cite{xhtml:oreilly}. Da moderne \Gls{webbrowser} unterschiedliche Technologien wie \gls{css} und \gls{javascript} zur
ergänzenden Darstellung von \gls{xhtml} nutzen, ist nicht sichergestellt, dass ein Zweig in einem \gls{xhtml}-Dokument einen geschlossenen Bereich in der
grafischen Darstellung des selben Dokumentes einnimmt. Umgekehrt können unterschiedliche Zweige aus dem \gls{dom}-Baum bei der Darstellung zu einem
geschlossenen Bereich zusammengeführt werden.

Als ein Bereich wird, im Rahmen dieser Arbeit, die Positionierung eines oder mehrerer Elemente innerhalb eines \gls{dom}-Baumes definiert, nicht die grafische
Darstellung in einem \gls{webbrowser}.

\section{Digitale Signaturen}
\index{Digitale Signaturen}
\label{sec:GrundlagenDefinitionen:DigitaleSignaturen}
Unter dem kryptographischen Sammelbegriff der digitalen Signatur versteht man eine zuverlässige und nicht fälschbare Kontrollinformation die über den Urheber
und mögliche Modifikationen des zu schützenden Inhalts Auskunft gibt \cite{kits}. Dabei kann eine digitale Signatur analog zu der Unterschrift einer Person
betrachtet werden, jedoch mit erweiterten Eigenschaften \cite{esig:bsi,angewandte-kryptographie,krypto-pki-internet}.

Der Begriff der digitalen Signatur wird als eine Information definiert, welche die Identität einer Urheber-Entität mit einer digitalen Information (Nachricht)
assoziiert und die drei kryptografischen Sicherheitsdienste Integrität, Authentizität und Verbindlichkeit erfüllt
\cite{hac,angewandte-kryptographie,krypto-pki-internet}.

\index{Digitale Signaturen!Integrität}
\label{sec:GrundlagenDefinitionen:DigitaleSignaturen:Integrität}
Durch \textbf{Integrität} wird garantiert, dass die Nachricht auf ihrem Weg zum Empfänger nicht verändert wurde \cite{niag}. Der Empfänger der Nachricht ist
damit in der Lage die Veränderung zu erkennen, kann jedoch nicht feststellen, welche Teile genau verändert wurden \cite{angewandte-kryptographie}.

\index{Digitale Signaturen!Authentizität}
\label{sec:GrundlagenDefinitionen:DigitaleSignaturen:Authentizität}
Der Empfänger einer Nachricht kann durch die \textbf{Authentizität} deren Urheber feststellen. Der Urheber der Nachricht weist dabei seine Identität seinem
Partner gegenüber durch ein Geheimnis oder eine Fertigkeit aus, die ihn für den Empfänger eindeutig kennzeichnet. Die Authentizität einer Nachricht setzt ihre
Integrität voraus, denn eine veränderte Nachricht ist auch nicht mehr authentisch \cite{kits,angewandte-kryptographie}.

\index{Digitale Signaturen!Verbindlichkeit}
\label{sec:GrundlagenDefinitionen:DigitaleSignaturen:Verbindlichkeit}
Als \textbf{Verbindlichkeit} einer Nachricht wird die Nachweisbarkeit des Urhebers einer Nachricht gegenüber Dritten definiert. Dieser Sicherheitsdienst schließt die
Eigenschaften der Authentizität und der Integrität mit ein \cite{kits,angewandte-kryptographie}.

\section{Digitale Signaturverfahren}
\index{Digitale Signaturverfahren}
\label{sec:GrundlagenDefinitionen:DigitaleSignaturen:Verfahren}
Ein digitales Signaturverfahren stellt ein Verfahren dar, mit welchem eine Entität ihre Identität an eine Information binden kann. Der Vorgang des Signierens
wird durch das Transformieren der Information und eines Geheimnisses, das nur dem Urheber gekannt ist, in eine Signatur, definiert \cite{hac}.

Ein Signaturverfahren bildet dabei für eine, zu signierende Nachricht, den passenden Hash-Wert und erzeugt mit dem privaten Schlüssel die Signaturinformation.
Diese Information wird nun in das, dem Verfahren entsprechende, Format transformiert und der Nachricht beigefügt.
Der Empfänger extrahiert die Signaturinformation aus der empfangenen Nachricht, bildet den Hash-Wert über die signierte Nachricht und vergleicht diesen mit dem
Hash-Wert aus der Signatur. Sind beide miteinander identisch, ist die Signatur gültig und die Integrität der Nachricht gewährleistet \cite{hac,singh:messages}.

Asymmetrische kryptographische Verfahren lassen sich zum Verschlüsseln wie auch zum Signieren von Informationen verwenden. Im Rahmen dieser Arbeit wird nur auf
den Signatur-Betrieb eingegangen.

\subsection{\glsentrytext{openpgp}}
\index{OpenPGP}
\label{sec:GrundlagenDefinitionen:DigitaleSignaturen:Verfahren:openpgp}
\gls{openpgp} wurde im Jahr 1991 von Philip Zimmermann \cite{zimmermann:pgp,pgp,intro-crypto} entwickelt, mit dem Ziel die digitale Kommunikation vor dem
Zugriff durch Dritte zu schützen. Eine der Motivationen während der Entwicklung waren die wiederholten Abhöraktionen der Bürgerrechtsbewegungen in den
Vereinigten Staaten von Amerika durch die dortigen Geheimdienste \cite{singh:messages}.

Der Name der aktuelle Spezifikation für das Datenformat ist \gls{openpgp} und seine Eigenschaften sind im RFC4880 festgelegt \cite{openpgp:ietf}. 
Jedes \gls{openpgp}-Zertifikat beinhaltet Informationen über:
\begin{itemize*}
    \item den öffentlichen Schlüssel,
    \item die Version des Datenformats,
    \item die Erstellungszeit,
    \item eine oder mehrere IDs des Inhabers des Zertifikats, meist E-Mail-Adressen,
    \item optionale Signaturen von anderen Benutzer, die das Zertifikat signiert haben,
    \item optionale Sub-Schlüssel, die ebenfalls dem Inhaber des Zertifikats gehören.
\end{itemize*}

\subsection{\glsentrytext{x509}}
\index{X.509v3}
\label{sec:GrundlagenDefinitionen:DigitaleSignaturen:Verfahren:x509}
Der \gls{itu} Standard \gls{x509} wurde 1996 veröffentlicht und beschreibt die Grundlage für Zertifikate und Widerrufslisten in der Syntax von \gls{asn1}. Die
Versionen 1 und 2 enthielten verschiedene Mängel, die einen Einsatz in produktiven Umgebungen erschwerten, so fehlten zum Beispiel Felder für wichtige
Informationen die erst mit Version 3 verfügbar gemacht wurden \cite{krypto-pki-internet}.

Jedes \gls{x509}-Zertifikat beinhaltet Informationen über:
\begin{itemize*}
    \item die Versionsnummer des Zertifikats,
    \item die Seriennummer,
    \item die \gls{oid} des Signaturverfahrens mit dem das Zertifikat signiert wurde,
    \item den Namen der \gls{ca} die das Zertifikat signiert hat,
    \item den Namen des Inhabers des Zertifikats,
    \item den öffentlichen Schlüssel,
    \item einen Gültigkeitszeitraum und
    \item den Verwendungszweck.
\end{itemize*}

Diese Informationen sind von dem übergeordneten Zertifikat, auch \gls{ca} genannt, digital signiert \cite{x509:ietf}.

\gls{x509} kommt zum Beispiel in Protokollen wie \gls{ssl}/\gls{tls}, \gls{ipsec} oder \gls{eap} zum Einsatz, da bei diesen eine streng hierarchische Struktur
der Vertrauensbeziehungen von Vorteil ist \cite{kits,tls:ietf}.

\section{Vertrauensbeziehungen}
\index{Vertrauensbeziehungen}
\label{sec:GrundlagenDefinitionen:Vertrauen}
Der Nutzen von digitalen Signaturen kommt erst zum Tragen, wenn auch sichergestellt wird, dass es eine Möglichkeit gibt, festzustellen, mit welcher Sicherheit
ein Schlüssel einer bestimmten Entität zugeordnet werden kann. Ist diese Sicherheit nicht gegeben, können keine Aussagen zur Authentizität einer Signatur
getätigt werden. Es müssen also Methoden existieren, um das eigene Vertrauen in einen Schlüssel auszudrücken und auch die Authentizität von Schlüsseln
zu überprüfen \cite{krypto-pki-internet}.

\gls{openpgp} und \gls{x509} verfolgen hierbei jeweils unterschiedliche Ansätze, um Vertrauensbeziehungen zwischen Schlüsseln zu definieren
\cite{krypto-pki-internet}.

\subsection{\glsentrydesc{wot}}
\index{Web Of Trust}
\label{sec:GrundlagenDefinitionen:Vertrauen:WOT}
Um die Ermittlung der Authentizität einer Signatur zu erleichtern, stellt \gls{openpgp} eine Funktionalität zur Beglaubigung von Schlüsseln zur Verfügung. Dabei
werden die Schlüssel von anderen \gls{openpgp}-Benutzern mit ihrer digitalen Signatur versehen. Das ermöglicht es, Rückschlüsse auf die Vertrauenswürdigkeit
von Schlüsseln zu ziehen, ohne von einer zentralen hierarchischen Struktur abhängig zu sein. Die Signatur auf einen Schlüssel wird üblicherweise nur nach der
Überprüfung der Identität des Schlüsselinhabers gewährt. Dies kann zum Beispiel durch ein Telefonat oder ein persönliches Treffen erfolgen, bei dem die
Prüfsumme des Schlüssels verglichen wird.

Wenn Alice den Schlüssel von Bob signiert, dann ermöglicht dies einer Person, die bereits dem Schlüssel von Alice vertraut, auch dem Schlüssel von Bob zu
vertrauen. Durch diese Delegation von Vertrauen bildet sich eine flache, vernetzte Struktur wie in \fref{fig:wot} dargestellt, die unter dem Namen \gls{wot}
bekannt ist. 

Jeder Benutzer von \gls{openpgp} ist somit in der Lage, sein Vertrauen in Schlüssel nach seinem eigenen Ermessen zu vergeben. Dafür stehen dem
Benutzer unterschiedliche Stufen des Vertrauens zur Verfügung die er mit seiner Signatur eines fremden Schlüssels verknüpfen kann. Sie reichen von "`Ich bin mir
nicht sicher"' bis zu "`Ich vertraue absolut"' \cite{pgp}.

\begin{figure}
\centering
\begin{tikzpicture}[node distance   = 4 cm]
    \tikzset{Key/.style = {shape          = rectangle,
                           rounded corners,
                           drop shadow,
                           fill           = orange,
                           draw           = black,
                           thick,
                           text           = black,
                           text centered,
                           inner sep      = 3pt,
                           outer sep      = 0pt,
                           minimum size   = 24 pt}}
    \tikzset{DirectTrust/.style   = {thick,
                                     double          = green,
                                     double distance = 1pt}}
    \tikzset{IndirectTrust/.style   = {thin,
                                       double          = blue,
                                       double distance = 1pt}}
    \tikzset{Label/.style =   {draw,
                               fill           = white,
                               text           = black}}
    \node[Key](A){Alice};
    \node[Key,right=of A](B){Bob};
    \node[Key,right=of B](C){Carmen};
    \node[Key,above=of B](D){Dorian};     
    \draw[DirectTrust](B) to node[Label]{vertraut} (D) ;
    \tikzset{DirectTrust/.append style = {bend left}}
    \draw[DirectTrust](A) to node[Label]{vertraut} (B);
    \draw[DirectTrust](B) to node[Label]{vertraut} (C);
    \draw[DirectTrust](D) to node[Label]{vertraut} (C);
    \tikzset{IndirectTrust/.append style = {bend left}}
    \draw[IndirectTrust](A) to node[Label]{vertraut indirekt} (D) ;
    \tikzset{IndirectTrust/.append style = {bend right}}
    \draw[IndirectTrust](A) to node[Label]{vertraut indirekt} (C) ;

\end{tikzpicture}

\caption{Schematische Darstellung des Web Of Trust}
\label{fig:wot}
\end{figure}

\subsection{\glsentrydesc{ht}}
\index{Hierachical Trust}
\label{sec:GrundlagenDefinitionen:Vertrauen:HT}
\gls{x509} bildet die Grundlage für ein hierarchisches System aus Zertifikaten, auch \gls{ht} genannt. Im Gegensatz zum \gls{wot}, welches einen Graphen der
beteiligten Schlüssel bildet, bestehen bei einem \gls{ht} eine oder mehrere baumartige Vertrauensstrukturen.
An der Spitze eines solchen Baumes steht eine Root-\gls{ca}, welche die Zertifikate der nächsten, im Baum tiefer gelegen Ebene signiert. Jeder Pfad, der von der
Wurzel, der Root-\gls{ca} des Baumes, zu einem der Zertifikate an einem Ast führt wird als \gls{cot} bezeichnet. \fref{fig:ht} stellt einen solchen Baum mit
mehreren \glspl{cot} dar.

Die Anzahl der Ebenen in einem \gls{ht}-Baum unterliegt keiner Beschränkung, jedoch muss einem Zertifikat explizit das Recht eingeräumt werden, andere
Zertifikate unter ihr zu signieren \cite{kits}. Im Internet existiert heute eine große Anzahl von \glspl{ca}, welche von Staaten, Unternehmen oder auch
Non-Profit-Organisationen betrieben werden. Zwischen diesen einzelnen Bäumen besteht keine Verbindung und auch keine Vertrauensbeziehung. Deshalb vertraut bei
einem \gls{ht} jeder Teilnehmer einer von ihm (oder seinem Softwaredistributor) ausgewählten Menge an \glspl{ca} \cite{ssliverse:eff,krypto-pki-internet}.

\begin{figure}
\centering
\begin{tikzpicture}[node distance   = 3 cm]
    \tikzset{Key/.style = {shape          = rectangle,
                           rounded corners,
                           drop shadow,
                           fill           = white,
                           draw           = black,
                           thick,
                           text           = black,
                           text centered,
                           inner sep      = 3pt,
                           outer sep      = 0pt,
                           minimum size   = 24 pt}}
    \tikzset{DirectTrust/.style   = {thick,
                                     double          = black,
                                     double distance = 1pt}}
    \tikzset{IndirectTrust/.style   = {thin,
                                       double          = white,
                                       double distance = 1pt}}
    \tikzset{Label/.style =   {draw,
                               fill           = white,
                               text           = black,
                               text centered,
                               text width     = 4em}}
    \node[Key](CA1){CA1 (Root)};
    \node[below=of CA1](AUX1){};
    \node[Key,left=of AUX1](CA1-1){CA1.1};
    \node[Key,right=of AUX1](CA1-2){CA1.2};
    \node[below=of CA1-1](AUX1-1){};
    \node[Key,left=of AUX1-1](Alice){Alice};
    \node[Key,right=of AUX1-1](Bob){Bob};
    \node[Key,below=of CA1-1](Carmen){Carmen};
    \node[Key,below=of CA1-2](Dorian){Dorian};
    \draw[DirectTrust,bend left](CA1-1) to node[Label]{vertraut} (CA1) ;
    \draw[DirectTrust,bend right](CA1-2) to node[Label]{vertraut} (CA1) ;
    \draw[DirectTrust,bend left](Alice) to node[Label]{vertraut} (CA1-1) ;
    \draw[DirectTrust,bend right](Bob) to node[Label]{vertraut} (CA1-1) ;
    \draw[DirectTrust](Carmen) to node[Label]{vertraut} (CA1-1) ;
    \draw[DirectTrust](Dorian) to node[Label]{vertraut} (CA1-2) ;
    \draw[IndirectTrust,bend right](Carmen) to node[Label]{vertraut indirekt} (Dorian) ;
    \draw[IndirectTrust,bend right](Alice) to node[Label]{vertraut indirekt} (Carmen) ;

\end{tikzpicture}

\caption{Schematische Darstellung eines \texorpdfstring{\protect\Glsentryname{ht}}{HT}}
\label{fig:ht}
\end{figure}


%%%%%%%%%%%%%%%%%%%%%%%%%%%%%%%%%%%%%%%
%%%     KAPITEL 3
%%%     Anforderungsnaalyse
%%%     2007/10/17
%%%%%%%%%%%%%%%%%%%%%%%%%%%%%%%%%%%%%%%
\chapter{Anforderungsanalyse f�r ein PPM-System}
\index{Anforderungsanalyse}%
\label{chp:Anforderungsanalyse}%


%%%------------------------------------
%%%---- Requirements Engineering -------------
%%%------------------------------------

\section{Requirements Engineering}
\index{Requirements Engineering}%
\label{sec:req:req}%


\subsection{Definition}
\index{Definition}
\label{sec:thg:req:definition}%

\subsection{Analysemodelle}
\index{Analysemodelle}
\label{sec:thg:req:modelle}%


%%%------------------------------------
%%%---- Prozessdefinition -------------
%%%------------------------------------

\section{Prozessdefinition}
\index{Prozessdefinition}%
\label{sec:req:process}%



%%%------------------------------------
%%%---- ANHANG ------------------------
%%%------------------------------------
\appendix

%%%------------------------------------
%%%---- NACHSPANN ---------------------
%%%------------------------------------
\backmatter %%% Bibliographie, Index

%%%------------------------------------
%%%     ABKUERZUNGSVERZEICHNIS
\clearpage%
\printglossaries
%%%------------------------------------
%%%     ABBILDUNGSVERZEICHNIS
\clearpage%
\listoffigures%
%%%------------------------------------
%%%     TABELLENVERZEICHNIS
\clearpage%
\listoftables%
%%%------------------------------------
%%%     LISTINGSVERZEICHNIS
\clearpage%
\lstlistoflistings%
%%%------------------------------------
%%%     LITERATURVERZEICHNIS
\clearpage%
\bibliographystyle{alphadin}%
\bibliography{misc/literatur}%
%%%------------------------------------
%%%     INDEX
\clearpage%
\printindex
\chapter*{}

\mbox{}

\textit{,,First they ignore you. Then they laught about you. Then
they fight you. And then you win.''}

\raggedleft{---  Mahatma Gandhi}
%

%%%------------------------------------
%%%     EO DOCUMENT
%%%------------------------------------
\end{document}%}}}

% vim:filetype=tex foldmethod=marker autoindent expandtab shiftwidth=4
% Local variables:
% mode: latex
% End:              
